% Generated by 'make docs' from the HELP_ strings in garlic-cli.cpp.
% Do not edit by hand.
\begingroup\sloppy
\begin{description}
\item[\texttt{--M} \textnormal{\textless{}int\textgreater{}}] The expected number of meioses since a recent common ancestor for \texttt{--weighted} calculation. \\ \textit{Default:} \texttt{7}
\item[\texttt{--auto-overlap-coef} \textnormal{\textless{}double1\textgreater{} ... \textless{}doubleN\textgreater{}}] \textless{}slope\textgreater{} \textless{}intercept\textgreater{}: coefficients of overlap\% = slope*log(density) + intercept, used by \texttt{--auto-overlap-frac}. The defaults (6.375 63.888) are an empirical fit to human SNP-array data. \\ \textit{Default:} \texttt{6.375 63.888}
\item[\texttt{--auto-overlap-frac} \textnormal{\textless{}bool\textgreater{}}] If set, GARLIC attempts to guess based on marker density. \\ \textit{Default:} \texttt{false}
\item[\texttt{--auto-winsize} \textnormal{\textless{}bool\textgreater{}}] If \texttt{--weighted} is set, guesses the best window size based on SNP density, otherwise initiates an ad hoc method for automatically selecting the \# of SNPs in which to calculate LOD scores. Starts at the value specified by \texttt{--winsize} and increases by \textless{}step size\textgreater{} SNPs until finished. \\ \textit{Default:} \texttt{false}
\item[\texttt{--auto-winsize-coef} \textnormal{\textless{}double1\textgreater{} ... \textless{}doubleN\textgreater{}}] \textless{}slope\textgreater{} \textless{}intercept\textgreater{}: coefficients of winsize = slope*log(density) + intercept, used by \texttt{--auto-winsize} with \texttt{--weighted}. The defaults (8.3235 138.0521) are an empirical fit to human SNP-array data; override them for other ascertainment schemes or species. \\ \textit{Default:} \texttt{8.3235 138.0521}
\item[\texttt{--auto-winsize-step} \textnormal{\textless{}int\textgreater{}}] Step size for automatic window selection algorithm. \\ \textit{Default:} \texttt{10}
\item[\texttt{--auto-winsize-threshold} \textnormal{\textless{}double\textgreater{}}] Smoothness criterion at which the \texttt{--auto-winsize} search stops. Lower values demand a smoother LOD score density and therefore a larger window. \\ \textit{Default:} \texttt{0.5}
\item[\texttt{--autosomes-only} \textnormal{\textless{}bool\textgreater{}}] Drop every chromosome that is not an autosome before calling: X, Y, Z, W and the mitochondrion, recognised by name, plus 23, 24, 25 and 26 when \texttt{--build} names a human assembly. A hemizygous genotype is written as a homozygous call and is indistinguishable from true autozygosity, so a hemizygous chromosome is otherwise called as one run spanning its whole length. Does not resolve bare 23-26 without a human \texttt{--build}: dropping them is right if they are PLINK's human codes and destroys real autosomes if they are not, so garlic refuses instead. See \texttt{--sex-system}. \\ \textit{Default:} \texttt{false}
\item[\texttt{--build} \textnormal{\textless{}string\textgreater{}}] Choose which genome build to use for centromere locations: hg18, hg19, hg38 or t2t-chm13 (T2T-CHM13v2.0, UCSC hs1). \\ \textit{Default:} \texttt{none}
\item[\texttt{--centromere} \textnormal{\textless{}string\textgreater{}}] Provide custom centromere boundaries. Format \textless{}chr\textgreater{} \textless{}start\textgreater{} \textless{}end\textgreater{}. \\ \textit{Default:} \texttt{none}
\item[\texttt{--chr} \textnormal{\textless{}string1\textgreater{} ... \textless{}stringN\textgreater{}}] Analyse only these chromosomes, e.g. \texttt{--chr} chr1 chr2 chrX. Names are matched after normalisation, so '1' and 'chr1' are the same. It is an error to name a chromosome that is not in the data. \\ \textit{Default:} \texttt{all chromosomes in the input}
\item[\texttt{--chr-lengths} \textnormal{\textless{}string\textgreater{}}] A file of chromosome lengths: two whitespace-separated columns, chromosome and length. A samtools faidx .fai is accepted unchanged -- its first two columns are already name and length and the rest is ignored -- so the index beside your reference can be passed as it is. The way to use \texttt{--froh-denominator} assembly on a non-human assembly, and on TPED input, where nothing in the data states a chromosome's length. Takes precedence over \texttt{--build} and over a VCF's \#\#contig header. \\ \textit{Default:} \texttt{none}
\item[\texttt{--cm} \textnormal{\textless{}bool\textgreater{}}] Measure ROH lengths in genetic distance units. This requires a mapfile. \\ \textit{Default:} \texttt{false}
\item[\texttt{--dump-docs} \textnormal{\textless{}string\textgreater{}}] Write the command line reference to stdout in a documentation format and exit: txt for the README block, tex for a LaTeX description list. Used by 'make docs' so the documentation cannot drift from the program. \\ \textit{Default:} \texttt{\_\_none}
\item[\texttt{--error} \textnormal{\textless{}double\textgreater{}}] The assumed genotyping error rate. \\ \textit{Default:} \texttt{-1}
\item[\texttt{--force} \textnormal{\textless{}bool\textgreater{}}] Overwrite existing output files. Without this, garlic refuses to clobber an existing \textless{}out\textgreater{}.roh.bed. \\ \textit{Default:} \texttt{false}
\item[\texttt{--freq-file} \textnormal{\textless{}string\textgreater{}}] A file specifying allele frequencies for each population for all variants. File format: CHR SNP POS ALLELE FREQ \textless{}chr\textgreater{} \textless{}locus ID\textgreater{} \textless{}allele\textgreater{} \textless{}freq\textgreater{} By default, this is calculated automatically from the provided data. \\ \textit{Default:} \texttt{none}
\item[\texttt{--freq-only} \textnormal{\textless{}bool\textgreater{}}] If set, calculates a freq file from provided data and then exits. Uses minimal RAM. \\ \textit{Default:} \texttt{false}
\item[\texttt{--froh} \textnormal{\textless{}bool\textgreater{}}] Also write \textless{}out\textgreater{}.froh.tsv: per-individual autozygous total and fraction, broken down by ROH size class. The denominator is stated in the file header. \\ \textit{Default:} \texttt{false}
\item[\texttt{--froh-denominator} \textnormal{\textless{}string\textgreater{}}] What \textless{}out\textgreater{}.froh.tsv divides the autozygous total by. analyzed (default): the span the markers cover -- from the first to the last locus of each analysed chromosome, less the centromere and any region excluded with \texttt{--par}. This is what a marker set can support on its own, and it moves with the markers, so FROH from two different arrays are not directly comparable. assembly: the FULL length of each analysed chromosome, nothing subtracted. This is what published FROH denominators are, and it is reproducible from an external table rather than from the data. Lengths come from \texttt{--chr-lengths}, from \texttt{--build}, or from a VCF's \#\#contig header, in that order; a chromosome with no length available stops the run rather than falling back. Not available with \texttt{--cm}: a base pair length is not a genetic length. \\ \textit{Default:} \texttt{analyzed}
\item[\texttt{--gl-type} \textnormal{\textless{}string\textgreater{}}] Which genotype-quality field to use.  With \texttt{--vcf} it is read from the FORMAT column; with \texttt{--tgls} it names the form of the single value per genotype in the .tgls file.  The two sources are NOT equivalent. With \texttt{--vcf} (recommended): GQ error = 10\textasciicircum{}(-GQ/10), the VCF spec's definition directly. PL The whole PL array is read and converted to a posterior:     P(g) proportional to 10\textasciicircum{}(-PL\_g/10),  error = 1 - P(called)/sum.   This is the only correct use of a VCF's PL.  A VCF normalises PL   so the CALLED genotype is exactly 0, so that one value carries no   information -- it is 0 for every call, confident or not.  The   information is in how much worse the alternatives are. GL The same, via PL = -10*GL.   A missing field is an error naming whichever of GQ/PL/GL the file   does carry, and a called genotype with '.' for the field is an   error rather than a silent fall back to \texttt{--error}. With \texttt{--tgls} (legacy; the format holds ONE value per genotype): GQ As above.  The only form a VCF's values can be dropped into. PL error = 1 - 10\textasciicircum{}(-PL/10), i.e. P(genotype CORRECT).   *** NOT the VCF PL field. *** Under this definition PL=0 means an   error of 0, and the integers map to 0, 0.21, 0.37, 0.50, ... --   none a realistic error rate (0.001 would be PL=0.00435).  Feeding   VCF PLs here is rejected.  Use \texttt{--vcf} \texttt{--gl-type} PL instead. GL error = 1 - 10\textasciicircum{}GL, with the same caveat. \\ \textit{Default:} \texttt{none}
\item[\texttt{--gmm-max-iter} \textnormal{\textless{}int\textgreater{}}] Maximum EM iterations when fitting ROH size classes. \\ \textit{Default:} \texttt{1000}
\item[\texttt{--gmm-tol} \textnormal{\textless{}double\textgreater{}}] Convergence tolerance for the ROH size-class EM fit. \\ \textit{Default:} \texttt{1e-05}
\item[\texttt{--haploid-chr} \textnormal{\textless{}string1\textgreater{} ... \textless{}stringN\textgreater{}}] Chromosomes haploid in every individual: mitochondrion, chloroplast. Dropped before calling, since a run of homozygosity on them is not a meaningful quantity. \\ \textit{Default:} \texttt{names matching M or MT, plus 26 under a human --build}
\item[\texttt{--help} \textnormal{\textless{}bool\textgreater{}}] Prints this help dialog. \\ \textit{Default:} \texttt{false}
\item[\texttt{--het-rate-bounds} \textnormal{\textless{}double1\textgreater{} ... \textless{}doubleN\textgreater{}}] \textless{}lo\textgreater{} \textless{}hi\textgreater{}: heterozygosity rates on the shared sex chromosome below which an individual reads as hemizygous and above which it reads as diploid. Between them it is ambiguous, and garlic neither calls ROH there for that individual nor counts its alleles. Used to check every recorded sex and to infer the ones that were not recorded. Fixed bounds rather than a two-component fit on purpose: a cohort that is entirely one sex has no bimodality to find, and a clustering rule would invent a split there. These fail the same way for every dataset. \\ \textit{Default:} \texttt{0.020000}
\item[\texttt{--kde-cut} \textnormal{\textless{}double\textgreater{}}] Extend the density grid this many bandwidths beyond the range of the data. \\ \textit{Default:} \texttt{3}
\item[\texttt{--kde-points} \textnormal{\textless{}int\textgreater{}}] Number of grid points at which the LOD score density is evaluated. Also sets the resolution at which the between-mode minimum (the LOD cutoff) is located. \\ \textit{Default:} \texttt{512}
\item[\texttt{--kde-subsample} \textnormal{\textless{}int\textgreater{}}] The number of individuals to randomly sample for LOD score KDE. If there are fewer individuals in the population all are used. The default, 0, uses every individual, which makes the selected LOD cutoff a deterministic function of the data alone. A positive value subsamples, which saves memory but makes the cutoff depend on \texttt{--seed}. \\ \textit{Default:} \texttt{0}
\item[\texttt{--kde-thin-step} \textnormal{\textless{}int\textgreater{}}] Take every Nth LOD score window when building the KDE. 0 (the default) uses the window size, which keeps the sampled windows non-overlapping; 1 uses every window. Replaces \texttt{--no-kde-thinning}, which is equivalent to 1 and is kept as a deprecated alias. \\ \textit{Default:} \texttt{0}
\item[\texttt{--ld-subsample} \textnormal{\textless{}int\textgreater{}}] The number of individuals to randomly sample for LD calculation during wLOD. If there are fewer individuals in the population all are used. Set \textless{}= 0 to use all individuals (will increase runtime). \\ \textit{Default:} \texttt{0}
\item[\texttt{--load-params} \textnormal{\textless{}string\textgreater{}}] Read flag values from a \textless{}out\textgreater{}.params.json written by a previous run. Flags given on the command line take precedence over the file, so a run can be repeated with one parameter changed. \\ \textit{Default:} \texttt{\_\_none}
\item[\texttt{--lod-cutoff} \textnormal{\textless{}double\textgreater{}}] For LOD based ROH calling, specify a single LOD score cutoff above which ROH are called in all populations.  By default, this is chosen automatically with KDE. \\ \textit{Default:} \texttt{-999999}
\item[\texttt{--map} \textnormal{\textless{}string\textgreater{}}] Provide a scaffold genetic map, sites that aren't present within this file are interpolated. Sites outside the bounds are filtered. This is required for wLOD calculations and any runs for which you wish to report ROH in units of cM. \\ \textit{Default:} \texttt{none}
\item[\texttt{--max-gap} \textnormal{\textless{}int\textgreater{}}] A LOD score window is not calculated if the gap (in bps) between two loci is greater than this value. \\ \textit{Default:} \texttt{200000}
\item[\texttt{--max-winsize} \textnormal{\textless{}int\textgreater{}}] Upper bound on the window size that \texttt{--auto-winsize} will try. The search previously had no bound and would grow past the number of loci if the smoothness criterion was never met. \\ \textit{Default:} \texttt{1000}
\item[\texttt{--mode-smooth-span} \textnormal{\textless{}int\textgreater{}}] Number of adjacent grid points used to smooth the density when locating its modes. Must be smaller than \texttt{--kde-points}. \\ \textit{Default:} \texttt{20}
\item[\texttt{--mu} \textnormal{\textless{}double\textgreater{}}] Mutation rate per bp per generation for \texttt{--weighted} calculation. \\ \textit{Default:} \texttt{1e-09}
\item[\texttt{--nclust} \textnormal{\textless{}int\textgreater{}}] Set number of clusters for GMM classification of ROH lengths. \\ \textit{Default:} \texttt{3}
\item[\texttt{--no-centromere} \textnormal{\textless{}bool\textgreater{}}] Treat every chromosome as having no assembled gap, instead of requiring \texttt{--build} or \texttt{--centromere}. For organisms with no centromere gap in the assembly, and for references garlic has no built-in table for (only hg18, hg19 and hg38 are built in; T2T-CHM13 is not). Mutually exclusive with \texttt{--build} and \texttt{--centromere}. \\ \textit{Default:} \texttt{false}
\item[\texttt{--no-kde-thinning} \textnormal{\textless{}bool\textgreater{}}] Set this flag to send all LOD score data to KDE function. This may dramatically increase runtime. \\ \textit{Default:} \texttt{false}
\item[\texttt{--out} \textnormal{\textless{}string\textgreater{}}] The base name for all output files. \\ \textit{Default:} \texttt{outfile}
\item[\texttt{--outdir} \textnormal{\textless{}string\textgreater{}}] Write all output files into this directory, which is created if it does not exist. Equivalent to prefixing \texttt{--out} with the path.
\item[\texttt{--overlap-frac} \textnormal{\textless{}double\textgreater{}}] The minimum fraction of overlapping windows above the LOD cutoff required to begin constructing a run. ROH will have a lower bound size threshold of WINSIZE*OVERLAP\_FRAC. If set to 0, GARLIC sets the value to the lowest sensible value: 1/winsize. \\ \textit{Default:} \texttt{0.25}
\item[\texttt{--par} \textnormal{\textless{}string1\textgreater{} ... \textless{}stringN\textgreater{}}] Pseudoautosomal regions on the shared sex chromosome, as \textless{}chr\textgreater{}:\textless{}start\textgreater{}-\textless{}end\textgreater{}, comma or space separated. They are diploid in both sexes and too short to call at a useful window size, so the loci inside them are dropped. ANY NUMBER of regions per chromosome: humans have two, other species have more, and a system with several shared sex chromosomes can have them on each. Coordinates are read in the assembly the input uses; there is deliberately no way to name a pseudoautosomal region by chromosome code, because PLINK's 25 is an ordinary autosome in most species. \texttt{--build} supplies the human X regions on its own, from the Genome Reference Consortium's definition for that assembly; see centromeres/par\_regions.txt. Naming any region here replaces them, and \texttt{--par} none keeps them all. \\ \textit{Default:} \texttt{the regions for --build, or none without one}
\item[\texttt{--par-file} \textnormal{\textless{}string\textgreater{}}] The same regions from a file, one per line, \textless{}chr\textgreater{} \textless{}start\textgreater{} \textless{}end\textgreater{}. '\#' comments and blank lines are ignored and .gz is accepted. Unlike \texttt{--centromere}, repeated rows for a chromosome ADD a region rather than replacing one. May be given together with \texttt{--par}; the regions accumulate. \\ \textit{Default:} \texttt{none}
\item[\texttt{--phased} \textnormal{\textless{}bool\textgreater{}}] Set if data are phased and you want to calculate r2 instead of hr2 while \texttt{--weighted} is set. Uses extra RAM. Has no effect on computations without \texttt{--weighted}. \\ \textit{Default:} \texttt{false}
\item[\texttt{--pool-populations} \textnormal{\textless{}bool\textgreater{}}] Analyse every individual together as one population, ignoring the population labels.  This is what garlic did before per-population analysis, and it restores the old output names as well: \textless{}out\textgreater{}.roh.bed rather than \textless{}out\textgreater{}.\textless{}POP\textgreater{}.roh.bed. Allele frequencies are then computed by pooling everyone, which inflates heterozygosity relative to any one population and biases the LOD scores for all of them.  Use it when the labels are not populations -- column 1 of a TFAM is the PLINK family ID, which is often one family or even one sample per value -- or to reproduce an earlier run. \\ \textit{Default:} \texttt{false}
\item[\texttt{--pop} \textnormal{\textless{}string\textgreater{}}] A file mapping sample ID to population, replacing the population labels that would otherwise come from the TFAM. Whitespace separated, '\#' comments and blank lines ignored, .gz accepted:   \textless{}sample\_id\textgreater{}  \textless{}population\textgreater{}  [sex] Two required columns in THAT order. Note this is the opposite order to a TFAM, whose first column is the population. The optional third column is sex in PLINK coding (1 male, 2 female, 0 or -9 unknown); unlike a TFAM, anything else is an error rather than being read as unknown. Rows are matched by sample ID, never by file order. Every sample in the data must have a row; extra rows are ignored and counted, so a cohort-wide file can be used against a subset. \\ \textit{Default:} \texttt{none}
\item[\texttt{--quiet} \textnormal{\textless{}bool\textgreater{}}] Suppress the progress bar and informational messages. Errors and warnings still go to stderr and to \textless{}out\textgreater{}.error. \\ \textit{Default:} \texttt{false}
\item[\texttt{--raw-lod} \textnormal{\textless{}bool\textgreater{}}] If set, LOD scores will be output to gzip compressed files. \\ \textit{Default:} \texttt{false}
\item[\texttt{--resample} \textnormal{\textless{}int\textgreater{}}] Number of resamples for estimating allele frequencies. When set to 0 (default), garlic will use allele frequencies as calculated from the data. \\ \textit{Default:} \texttt{0}
\item[\texttt{--seed} \textnormal{\textless{}int\textgreater{}}] Random seed, for reproducible runs. Affects \texttt{--kde-subsample}, \texttt{--ld-subsample} and \texttt{--resample}. If 0 (default), a nondeterministic seed is drawn and written to the log file so the run can be repeated exactly. \\ \textit{Default:} \texttt{0}
\item[\texttt{--sex-chr} \textnormal{\textless{}string1\textgreater{} ... \textless{}stringN\textgreater{}}] The chromosome diploid in the homogametic sex and hemizygous in the heterogametic one: X under xy, Z under zw. Only needed when the name is not the conventional one, e.g. \texttt{--sex-system} zw \texttt{--sex-chr} LG12. Requires \texttt{--sex-system}. \\ \textit{Default:} \texttt{the conventional name for --sex-system}
\item[\texttt{--sex-chr-degenerate} \textnormal{\textless{}string1\textgreater{} ... \textless{}stringN\textgreater{}}] The chromosome carried only by the heterogametic sex: Y under xy, W under zw. No individual can carry a run of homozygosity on it -- hemizygous in one sex and absent in the other -- so it is dropped before calling. \\ \textit{Default:} \texttt{the conventional name for --sex-system}
\item[\texttt{--sex-system} \textnormal{\textless{}string\textgreater{}}] Which sex is heterogametic, and therefore which individuals are hemizygous on the shared sex chromosome:   xy    individuals coded 1 (male) in the TFAM or \texttt{--pop} are heterogametic   zw    individuals coded 2 (female) are heterogametic   none  assert that every chromosome is diploid in every individual, and         stop reporting sex chromosomes entirely PLINK sex coding is 1 male and 2 female whatever the species' sex determination system is, so in a ZW species the HOMOGAMETIC sex is coded 1. This is the one thing no amount of genotype data can supply, which is why it is a flag. It is not an analysis path: it fills in a per-chromosome, per-individual expected ploidy that one code path then reads. Leaving it unset is NOT the same as 'none'. If a sex chromosome is present and this is unset, garlic stops rather than analysing it as an autosome. \\ \textit{Default:} \texttt{unset}
\item[\texttt{--sexchr-lod-cutoff} \textnormal{\textless{}double\textgreater{}}] Call the shared sex chromosome at this LOD score cutoff instead of the one estimated on the autosomes. Every run logs what the sex chromosome's own windows would have given, which is where the value for this comes from. The autosomal cutoff is the default because one chromosome's windows from part of a cohort often cannot locate a second minimum in the density at all; that is also why this is a value you supply rather than a switch that asks garlic to re-estimate. \\ \textit{Default:} \texttt{the autosomal cutoff}
\item[\texttt{--sexchr-winsize} \textnormal{\textless{}int\textgreater{}}] Build windows of this many SNPs on the shared sex chromosome instead of the size used on the autosomes. The overlap FRACTION is unchanged, so the number of windows a SNP must fall in scales with the window. A LOD score cutoff belongs to the window size it was estimated at, so with this the autosomal cutoff no longer applies to the sex chromosome; give \texttt{--sexchr-lod-cutoff} as well. \\ \textit{Default:} \texttt{the autosomal window size}
\item[\texttt{--size-bounds} \textnormal{\textless{}double1\textgreater{} ... \textless{}doubleN\textgreater{}}] Specify the size class boundaries ROH boundaries.  By default, this is chosen automatically with a 3-component GMM.  Must provide numbers in increasing order. \\ \textit{Default:} \texttt{-1.000000}
\item[\texttt{--tfam} \textnormal{\textless{}string\textgreater{}}] A tfam formatted file containing population and individual IDs. \\ \textit{Default:} \texttt{none}
\item[\texttt{--tgls} \textnormal{\textless{}string\textgreater{}}] A tgls file containing one per-genotype likelihood value per individual, in the same row order as the tped.  See \texttt{--gl-type} for what the values must mean: for PL and GL this is NOT the VCF convention. \\ \textit{Default:} \texttt{none}
\item[\texttt{--threads} \textnormal{\textless{}int\textgreater{}}] The number of threads to use. Applies to LOD score calculation, the LD calculations under \texttt{--weighted}, the LOD score KDE, and ROH assembly. \\ \textit{Default:} \texttt{1}
\item[\texttt{--tped} \textnormal{\textless{}string\textgreater{}}] A tped formatted file containing map and genotype information. \\ \textit{Default:} \texttt{none}
\item[\texttt{--tped-missing} \textnormal{\textless{}char\textgreater{}}] Single character missing data code for TPED files. \\ \textit{Default:} \texttt{0}
\item[\texttt{--vcf} \textnormal{\textless{}string\textgreater{}}] A VCF file (plain or gzipped) to read genotypes from, instead of \texttt{--tped}. Biallelic SNVs with a GT field are used; multiallelic sites, indels and symbolic ALTs are skipped and the counts reported. Sample names come from the \#CHROM line, so \texttt{--tfam} is neither needed nor accepted; use \texttt{--pop} to supply population labels. Sites must be grouped by chromosome. Ploidy must be 2, except on a chromosome declared with \texttt{--sex-system} or \texttt{--sex-chr}, where a haploid call is read as hemizygous. With \texttt{--phased} every call must use '|'. \\ \textit{Default:} \texttt{none}
\item[\texttt{--vcf-pass-only} \textnormal{\textless{}bool\textgreater{}}] Skip sites whose FILTER column is neither PASS nor '.'. Without this, such sites are KEPT and their number is reported, because a FILTER column reflects the caller's thresholds rather than garlic's. \\ \textit{Default:} \texttt{false}
\item[\texttt{--verbose} \textnormal{\textless{}bool\textgreater{}}] Show the progress bar even when stderr is not a terminal. By default it is drawn only on a terminal, because it works by emitting backspaces. \\ \textit{Default:} \texttt{false}
\item[\texttt{--version} \textnormal{\textless{}bool\textgreater{}}] Print the version and exit. \\ \textit{Default:} \texttt{false}
\item[\texttt{--weighted} \textnormal{\textless{}bool\textgreater{}}] Compute LOD scores weighted by LD and probability of mutation. \\ \textit{Default:} \texttt{false}
\item[\texttt{--winsize} \textnormal{\textless{}int\textgreater{}}] The window size in \# of SNPs in which to calculate LOD scores. \\ \textit{Default:} \texttt{0}
\item[\texttt{--winsize-multi} \textnormal{\textless{}int1\textgreater{} ... \textless{}intN\textgreater{}}] Provide several window sizes (in \# of SNPs) to calculate LOD scores. LOD score KDEs for each window size will be output for inspection. \\ \textit{Default:} \texttt{-1}
\item[\texttt{-h} \textnormal{\textless{}bool\textgreater{}}] Prints this help dialog. \\ \textit{Default:} \texttt{false}
\end{description}
\endgroup
